\documentclass[12pt]{article}

\usepackage[utf8]{inputenc}
\usepackage{amsmath, amssymb, amsthm}
\usepackage{geometry}
\geometry{a4paper, margin=2.5cm}

\title{Teorema de Extensie (Carathéodory)}
\author{}
\date{}

\theoremstyle{plain}
\newtheorem{theorem}{Teoremă}

\begin{document}

\maketitle

\section{Definiție: Pre-măsură}

Fie $(X, \mathcal{A})$ un spațiu măsurabil, unde $\mathcal{A}$ este o
algebră de mulțimi.

O aplicație


\[
\mu_0 : \mathcal{A} \to [0, \infty]
\]


se numește \textbf{pre-măsură} dacă:

\begin{enumerate}
    \item $\mu_0(\emptyset) = 0$;
    \item este \textbf{sigma-aditivă} pe $\mathcal{A}$, adică pentru orice familie
    de mulțimi disjuncte $A_1, A_2, \dots \in \mathcal{A}$ cu
    $\bigcup_{n=1}^\infty A_n \in \mathcal{A}$,
    

\[
        \mu_0\!\left(\bigcup_{n=1}^\infty A_n\right)
        = \sum_{n=1}^\infty \mu_0(A_n).
    \]


\end{enumerate}

O pre-măsură este o măsură ,,incompletă'', definită doar pe o algebră.

\section{Teorema de Extensie (Carathéodory)}

\begin{theorem}[Teorema de Extensie]
Fie $\mathcal{A}$ o algebră de mulțimi pe $X$ și fie
$\mu_0 : \mathcal{A} \to [0,\infty]$ o pre-măsură.

Atunci există o \textbf{sigma-algebră} $\sigma(\mathcal{A})$ și o
\textbf{măsură} $\mu : \sigma(\mathcal{A}) \to [0,\infty]$ astfel încât:



\[
\mu(A) = \mu_0(A), \qquad \forall A \in \mathcal{A}.
\]



Mai mult, dacă $\mu_0$ este \textbf{sigma-finită}, atunci această măsură
$\mu$ este \textbf{unică}.
\end{theorem}

\section{Interpretare}

Teorema afirmă că:

\begin{itemize}
    \item dacă definim o măsură pe o algebră simplă de mulțimi,
    \item atunci ea poate fi extinsă riguros la întreaga sigma-algebră generată,
    \item iar dacă pre-măsura este sigma-finită, extensia este unică.
\end{itemize}

Aceasta este teorema fundamentală care permite construirea măsurii Lebesgue
și a multor alte măsuri importante în analiză și probabilități.

\section{Exemplu}

Pe $\mathbb{R}$, definim pe algebră mulțimilor de tip interval finit:



\[
\mu_0((a,b]) = b - a.
\]



Aceasta este o pre-măsură.  
Prin teorema de extensie, ea se extinde în mod unic la măsura Lebesgue
pe $\mathcal{B}(\mathbb{R})$.

\end{document}
